\par\clearpage
\noindent\null\hfill{\small\textbf{ZUR HANNOVERSCHEN TRIANGULATION.}}\hfill\makebox[0pt][r]{\fontsize{13}{16}\selectfont 411}\par\medskip
\noindent sich ferner die Möglichkeit, den Falkenberg mit Wilsede vermittelst eines
Durchhaus durch das Becklinger Holz zu verbinden, und die Hoffnung, dass
die beiden Punkte Hauselberg und Garssen vermittelst eines oder einiger
Zwischenpunkte und einiger Durchhaue durch die Waldungen, welche sie
scheiden, würden verknüpft werden können. Allein rücksichtlich solcher
Durchhaue, ohne welche in diesen Gegenden schlechterdings nicht durchzukommen
ist, glaubte ich mir zwei Gesetze auflegen zu müssen; erstlich, so
viel irgend möglich, allen edlern Holzarten (als Eichenwäldern) auszuweichen,
und zweitens die Richtung der Durchhaue \textls[120]{vorher} mit \textls[120]{äusserster Präcision}
zu bestimmen, so dass sie so schmal wie möglich ausfallen sollten, und
auch nicht Ein Stamm ohne Noth gefällt zu werden brauchte. Um aber
dies zu erreichen, mussten schon sehr scharfe Messungen, die zum Theil
künstlich arrangirt und combinirt werden mussten, vorangehen, und diese
liessen sich mit Genauigkeit und mit dem möglich geringsten Zeitaufwand erst
dann ausführen, wenn die Hauptmessungen erst selbst bis in diese Gegend
vorgerückt waren, und ich darf hier im voraus bemerken, dass ich diese
Zwecke späterhin auch zu meiner vollkommensten Zufriedenheit erreicht habe.

Die Zwischenzeit von der Rückkehr von der Recognoscirungsreise bis
zum eigentlichen Anfang der Messungsoperationen verwandte ich dazu, die
Instrumente ganz in gebrauchfertigen Stand zu setzen, alles nöthige vorzubereiten,
und dem Hauptmann \textsc{Müller} und dem Lieutenant \textsc{Hartmann} diejenigen
Instructionen zu geben, die zu einem vollkommenen Ineinandergreifen
der Operationen erforderlich waren. Am 16.~Junius trat ich die Reise nach
Lichtenberg, dem ersten in diesem Jahre vorzunehmenden Dreieckspunkte an,
wohin mein Sohn mit einem Theile der Instrumente schon vorausgereist war.

Die wirklichen Messungsarbeiten haben in diesem Jahre vier Monate gedauert,
nemlich bis Mitte Octobers, und während dieser Zeit sind \textls[120]{neun}
Hauptdreieckspunkte absolvirt: im Jahr 1821 konnten in einer nicht viel
kürzern Zeit nur fünf vorgenommen werden; dies so viel raschere Fortschreiten
ist hauptsächlich der Vermehrung der Zahl der Gehülfen und der
Heliotrope zuzuschreiben; allein noch wichtiger ist der daraus erhaltene Gewinn
in der Vergrösserung der Genauigkeit der Messungen selbst.\quad\mbox{.\ .\ .\ .\ .\ .}

Der Zweck der Triangulirung, als Theil der Gradmessung, ist, die Göttinger
Sternwarte durch ein zusammenhängendes System von Dreiecken mit den \mbox{däni-}
\par\smallskip\noindent\hfill\makebox[0.18\linewidth][l]{52\textsuperscript{*}}\par
\clearpage
\noindent\makebox[0pt][l]{\fontsize{13}{16}\selectfont 412}\hfill{\small\textbf{NACHLASS.}}\hfill\null\par\medskip
\noindent schen Dreiecken zu verbinden, und dazu war es am vortheilhaftesten, die Dreiecke
so gross wie möglich einzurichten, und die Dreieckspunkte im allgemeinen
auf den höchsten Stellen, die die weiteste Aussicht gewähren, zu wählen.
Diese Punkte haben an und für sich grösstentheils kein unmittelbares Interesse
für die Geographie des Königreiches. Ich habe aber überall, neben dem
Hauptzwecke, meine Operationen für diese nach Möglichkeit nützlich zu
machen gesucht. Die Lage aller Örter, die von mehr als einem Hauptdreieckspunkte
sichtbar waren, habe ich sorgfältig bestimmt, manche mit einer
Genauigkeit, die der in der Lage der Hauptdreieckspunkte gleich kommt. Ich
nenne davon die Städte Hannover, Braunschweig, Celle, Lüneburg, Neustadt
am Rübenberge, Burgdorf. Die Anzahl der Dörfer, deren Lage ich bestimmt
habe, ist sehr gross. Die Bahn ist gebrochen, diese Erndte über einen
grössern Theil des Königreichs, oder über das ganze, auszudehnen.\quad\mbox{.\ .\ .\ .\ .\ .}

\bigskip
\begin{center}
[Aus dem Bericht vom 16.~Februar 1825 ȟber die im Jahre 1824\par
ausgeführten trigonometrischen Arbeiten.«]
\end{center}

\par\noindent\hspace*{\parindent}\mbox{.\ .\ .\ .\ .\ .}\quad Resultate: Der Bericht und die beigefügte Karte zeigen, dass
die Arbeiten des Jahres 1824 mehrfache Übergänge von den Dreiecken der
frühern Jahre bis Bremen darbieten; der einfachste [die Dreiecke Wilsede-Falkenberg-Elmhorst,
Wilsede-Elmhorst-Litberg, Wilsede-Litberg-Hamburg,
Wilsede-Litberg-Zeven, Wilsede-Zeven-Steinberg und Zeven-Steinberg-Bremen
umfassend\,*)] ist mit starken vollen Linien gezeichnet, und zwei neue Dreiecke
[Zeven-Bremen-Brillit und Brillit-Bremen-Garlste] sind noch an die Seite
Bremen-Zeven angeknüpft. Wäre es möglich gewesen, jenes einfachste System
gleich anfangs ausfindig zu machen, so hätten allerdings die andern mit
schwachen vollen Linien gezeichneten Dreiecke [bei denen die Punkte Bullerberg,
Bottel und Brüttendorf Eckpunkte sind] ganz wegfallen können. Allein
die Berichterstattung zeigt, nach wie vielen Schwierigkeiten der Plan zu jenem
erst ausgemittelt werden konnte, und die präcise schnelle Ausführung der
verschiedenen dazu erforderlichen Durchhaue wäre gleichfalls ohne vorgängige
\par\vfill\noindent\rule{0.25\linewidth}{0.4pt}\par\smallskip
{\footnotesize\noindent\hspace*{\parindent}[*) Siehe die Dreiecksskizze auf S.~299].\par}
\clearpage
\noindent\null\hfill{\small\textbf{ZUR HANNOVERSCHEN TRIANGULATION.}}\hfill\makebox[0pt][r]{\fontsize{13}{16}\selectfont 413}\par\medskip
\noindent schon sehr genaue Kenntniss der Lage der Plätze ganz unthunlich gewesen.
Bei dem heutigen mathematischen Zustande der höhern Geodäsie dürfen übrigens
auch die letztern Dreiecke, die schwach gezeichneten, keinesweges als
überflüssig betrachtet werden: vielmehr muss ihre nach ganz bestimmten Principien
anzustellende Berücksichtigung mit dazu beitragen, die Schärfe der Endresultate
zu vergrössern.\quad\mbox{.\ .\ .\ .\ .\ .}

Endlich ist es auch noch von grosser Wichtigkeit, dass durch die drei
Dreiecke zwischen den fünf Punkten Falkenberg, Elmhorst, Wilsede, Litberg,
Hamburg ein neuer Übergang von den südlichen Dreiecken im Königreich
Hannover bis Hamburg erreicht worden ist, welcher vor dem frühern um
vieles complicirtern [über Hauselberg, Wulfsode, Timpenberg, Nindorf und
Lüneburg] vorzuziehen ist, und daher nach den vorhin angedeuteten Grundsätzen
die Genauigkeit der Resultate verdoppeln wird.\quad\mbox{.\ .\ .\ .\ .\ .}

\bigskip
\begin{center}
[Aus dem Bericht vom 21.~November 1827, »betreffend die weitere Ausdehnung\par
der Gradmessungsarbeiten.«]
\end{center}

\par\noindent\hspace*{\parindent}\mbox{.\ .\ .\ .\ .\ .}\quad Es wurden 1821---1823 bei Messung der Dreieckskette bis Hamburg
verausgabt 11000 Thaler, und 1824, 1825 für die von da westlich bis Ostfriesland
geführte Dreieckskette etwa 7000 Thaler. Von diesen Kosten ist aber
abzurechnen, was wegen Anschaffung von Instrumenten und wegen Abholens
des englischen Zenithsectors von Altona nach Göttingen verausgabt ist, und
zwischen 2500 und 3000 Thalern betragen haben mag, so dass die eigentlichen
Triangulirungskosten etwa 15000 Thaler betragen haben mögen. Nun scheint
nach der Übersichtskarte der Inbegriff der noch nicht berührten Landestheile
wohl nicht viel grösser zu sein, als die mit Dreiecken bereits überzogene
Fläche, und bei aller Ungewissheit, in der ich wegen der Schwierigkeiten des
Terrains bin, ist es doch kaum wahrscheinlich, dass sie grösser sein können,
als diejenigen, womit ich besonders 1822 und 1824 zu kämpfen gehabt habe.
Wenn ich nun ausserdem bemerke, dass die Operationen, deren Hauptzweck
die Vervollkommnung der Landes-Geographie ist, auch bei einer würdigen
Ausführung doch nicht \textls[120]{den} Grad von äusserster Schärfe der Messungen erfordern,
welcher bei einer eigentlichen Gradmessung verlangt wird, so scheint
\par\clearpage
\noindent\makebox[0pt][l]{\fontsize{13}{16}\selectfont 414}\hfill{\small\textbf{NACHLASS.}}\hfill\null\par\medskip
\noindent die Hoffnung nicht ungegründet, dass die Erweiterung der Triangulirung über
die noch nicht berührten Theile des Königreichs sich mit einer geringern
Summe und vielleicht mit 12000 Thalern bestreiten lassen werde\,*).

Um nun aber eine solche Triangulirung für die Vervollkommnung der
Geographie möglichst nützlich zu machen, wird man sich nicht darauf einschränken
müssen, bloss Netze von Hauptdreiecken der ersten Ordnung über
die betreffenden Landestheile auszuführen, sondern damit die Bestimmung der
Lage einer möglichst grossen Anzahl secundärer Punkte verbinden, namentlich
solcher, die scharfe Bestimmungen zulassen, und in der Regel Jahrhunderte
dauern, also besonders der Kirchthürme. Ich habe mir diese Rücksicht
schon bei den frühern Messungen zur Pflicht gemacht, obwohl sie dem Hauptzweck
untergeordnet bleiben musste, und die Anzahl der bei jenen Messungen
bestimmten Punkte beträgt schon über 500.\quad\mbox{.\ .\ .\ .\ .\ .}

Diese Angabe der Lage einer grossen Anzahl fester Punkte in Zahlen
(wie viel nemlich nördlich oder südlich, westlich oder östlich, von einem beliebigen
Anfangspunkte, z.~B. der Göttinger Sternwarte), \textls[120]{bis auf wenige
Fuss genau}, muss als die Hauptausbeute der Operationen in topographischer
Rücksicht betrachtet werden. Sie behält auf Jahrhunderte einen bleibenden
Werth, insofern die Mehrzahl der Punkte bleibt, wenn auch im Laufe der
Zeit einige untergehen, und die dadurch etwa entstehenden Veränderungen
sind leicht zu ergänzen. Sie bildet eine sichere Grundlage für alle Detailaufnahmen:
alle die Unsicherheiten, welche Aufnahmen ohne solche feste Anhaltspunkte
erschweren, entstellen und ihre Vereinigung zu einem fehlerfreien
Ganzen unmöglich machen, fallen dabei ganz weg; nachlässige Arbeiter erhalten
dadurch eine strenge unausweichliche Controlle; jede Messtischplatte
wird unabhängig von der andern bearbeitet, kein Fehler pflanzt sich also auf
andere Blätter fort; endlich vereinigen sich alle einzelnen Blätter von selbst
zu einem genau orientirten und überall zusammenpassenden Ganzen. Es ist
einleuchtend, dass die grossen Kosten, welche Detailaufnahmen von bedeutendem
Umfange allezeit machen, durch einen solchen sichern Gang in einem hohen
Grade vermindert werden müssen; aber dieser sichere Gang ist es nicht allein,
\par\vfill\noindent\rule{0.25\linewidth}{0.4pt}\par\smallskip
{\footnotesize\noindent\hspace*{\parindent}*) Die im laufenden Jahr 1827 erforderlichen Kosten kommen hiebei nicht in Betracht, da ihr Gegenstand zu dem rein astronomischen Theile der Gradmessung gehört.\par}
\clearpage
\noindent\null\hfill{\small\textbf{ZUR HANNOVERSCHEN TRIANGULATION.}}\hfill\makebox[0pt][r]{\fontsize{13}{16}\selectfont 415}\par\medskip
\noindent was die Arbeit beschleunigt; sehr wichtig ist in dieser Beziehung auch der
Umstand, dass der Gebrauch der Messkette dadurch fast ganz überflüssig und
nur ausnahmsweise nöthig wird, da die Triangulirung die Grundlinien schon
von selbst gibt, und mit einer Schärfe, welche die gewöhnliche Kette gar
nicht einmal geben könnte.

Allein auch, wo schon Detailaufnahmen vorhanden sind, wie bei den
meisten Ämtern des frühern Bestandes des Königreiches, bieten die festen
Punkte das Mittel dar, die aus der Zusammensetzung entstandenen Fehler zu
berichtigen, und dadurch selbst Karten, die sich auf unvollkommene Aufnahme-Methoden
gründen, wenn sie sonst im kleinen Detail gut sind, zu Darstellungen
umzuarbeiten, die auch höhern Anforderungen Genüge leisten können.

Was demnach die Maassregeln betrifft, um die Triangulirungen zur Vervollkommnung
der Geographie möglichst nützlich zu machen, so sind dabei
die bereits ausgeführten Messungen von den eventuell über andere Landestheile
künftig zu erstreckenden zu unterscheiden.

Bei letztern wird die Gewinnung genauer Bestimmung einer möglichst
grossen Anzahl fester Punkte gleich als Hauptzweck berücksichtigt werden
müssen.

Bei den bereits ausgeführten Messungen hingegen ist allerdings diese
Rücksicht nur als untergeordnet betrachtet gewesen; allein da ich, wie schon
erwähnt, dieselbe doch stets im Auge gehabt habe, so viel, ohne das Hauptgeschäft
zu hemmen, geschehen konnte, so müssen hinsichtlich des Erfolges
hier abermals die nördlichen Gegenden von den südlichen unterschieden werden.

In der nördlichen (grössern) Hälfte, d.~i. etwa von der Stadt Hildesheim
an bis zum Meere, also in dem flachen Theile des Landes, ist die Ausbeute in
der erwähnten Beziehung so ergiebig gewesen, dass wenig oder nichts zu
wünschen übrig bleibt.\quad\mbox{.\ .\ .\ .\ .\ .}

In dem südlichsten Theil des Königreichs hingegen ist die Anzahl der
scharf bestimmten Kirchthürme viel kleiner, da theils wegen der Grösse der
Dreiecke, theils wegen der gebirgigen Beschaffenheit des Landes nur wenige
Thürme von mehr als Einem Hauptdreieckspunkte aus zugleich sichtbar waren.
Für die Vervollkommnung der Geographie des Königreichs, und namentlich
um einer Detailaufnahme der südlichen Theile des Hildesheimschen und des
Eichsfeldes ähnliche sichere Grundlagen zu verschaffen, würde es daher \mbox{aller-}
\par
